\documentclass[11pt]{article}

% =====================================================================
% Internal APIs Were Not All You Needed
% The signed sovereign computation stack — Unbrowse, ZK identity, FDRY.
% A real whitepaper that knows it is a whitepaper. Every technical claim
% carries a real citation; claims are tagged [implemented] or [proposed]
% so nobody gets sold a vibe with a footnote.
% =====================================================================

\usepackage[margin=1in]{geometry}
\usepackage{hyperref}
\usepackage{amsmath}
\usepackage{amssymb}
\usepackage{enumitem}
\usepackage{xcolor}
\hypersetup{colorlinks=true, linkcolor=blue, urlcolor=blue, citecolor=blue}

\newcommand{\impl}{\textcolor{green!55!black}{\small[shipped]}}
\newcommand{\prop}{\textcolor{orange!80!black}{\small[proposed]}}
\newcommand{\CA}{\texttt{2ZiSPGncrkwWa6GBZB4EDtsfq7HEWwkwsPFzEXieXjNL}}

\title{\textbf{Internal APIs Were Not All You Needed}\\
\large A Signed, Layer-Descending Stack for Agent Computation:\\
Identity, Authentication, and Zero-Knowledge Privacy across the Stack}

\author{Lewis Tham\\ Unbrowse AI \\ \texttt{lewis@unbrowse.ai}}
\date{March 2026}

\begin{document}
\maketitle

\begin{abstract}
The field agreed, sensibly, that to let an agent use the web you mainly need a
good search API and a tidy way to call internal endpoints. This is true. It is
also, with apologies to a more famous paper,\footnote{The title is an affectionate
mugging of ``Attention Is All You Need.'' Yes, we know. That \emph{is} the bit.}
\emph{not all you needed.} An agent stops being a polite HTTP client the instant
it has to click something, run a shell command, open a real browser, or put bytes
on a socket that a bot-detector is staring at. The interesting work does not
politely confine itself to the layer where JSON comes out. This paper makes one
claim, and tries not to oversell it: the discipline that makes a \emph{single}
signed request trustworthy --- sign it, seal it, cache it, prove who you are
without showing your hand --- is the same discipline, unchanged, all the way down
the stack (screen $\to$ browser $\to$ CLI $\to$ OS $\to$ kernel $\to$ packet).
Every layer is signed by one Ed25519 wallet key; credentials are bound to that
identity by zero-knowledge proof and revealed only under signature; every
intermediate result is content-addressed and cached, with the signed commitment
appended to a hash-chained ledger (value off-chain, root on-chain); and nothing
leaves any layer unsigned. The maintenance economy that keeps those routes fresh
--- bonded accountability and proof of indexing --- is the concern of the companion
Maintenance Network paper; this paper is about the security, authentication, and
privacy of the stack itself. We cite a real primitive for
each layer instead of inventing one, and we label what is \textbf{[shipped]}
versus \textbf{[proposed]} so the reader always knows whether they are looking at
a product or a promise. The short version: internal APIs were a great first
layer. They were never the last one.
\end{abstract}

\section{Introduction: one layer is not the stack}

The agent-tooling world converged on a genuinely good wedge: stop paying a browser
to re-derive the web on every call, and instead look up a reusable route and just
execute it~\cite{fdry}. Unbrowse does exactly that --- intent $\to$ ranked
endpoint shortlist $\to$ execute --- and it wins on cost, latency, and
reliability. We are not here to relitigate that. It works. \impl{}

We are here about the quiet assumption underneath it: that the load-bearing part
of agent work lives at the HTTP layer, where things are clean. They are not clean.
A login form is a screen-layer fact. A session cookie is a browser-layer fact. A
keychain entry is an OS-layer fact. A TLS fingerprint is a packet-layer fact that
a defender is fingerprinting whether or not you were thinking about
it~\cite{ja3,ja4}. An agent that is actually \emph{sovereign} over its own actions
has to be the same coherent entity across all of those layers --- not just the one
that happens to return \texttt{application/json}.

So the thesis fits on one line: \emph{the layering is the point, and the same
covenant holds at every layer.} The end-to-end argument in system
design~\cite{e2e} already told us where each function belongs; we are simply
taking it at its word and running it down a signed agent stack to see if the shape
survives. (It does. That is the paper.)


\section{Related work}

Our contribution is not a new agent or a new chain; it is the claim that one signed
discipline holds across every layer an agent touches, and that the route economy is
its natural settlement. We situate that against four lines of prior work.

\paragraph{Web agents and the browser-first assumption.} A large body of work has
agents drive real browsers: ReAct interleaves reasoning and action~\cite{react},
and benchmarks such as Mind2Web~\cite{mind2web}, WebArena~\cite{webarena}, and the
GPT-4V grounding study SeeAct~\cite{seeact} measure agents on live or simulated
sites. This work treats the rendered page as the interface. We take the opposite
stance the Maintenance Network paper argues~\cite{fdry}: the browser is a fallback,
and the durable interface is the first-party API the page already calls. Our stack
keeps the browser as the bottom of a descent, not the default.

\paragraph{Tool and API invocation.} Toolformer teaches a model \emph{when} to call
an API~\cite{toolformer}; Gorilla teaches it to emit \emph{correct} calls across
thousands of APIs~\cite{gorilla}; the Model Context Protocol standardises how a model
connects to tools and data~\cite{mcp}. These solve single-agent tool use. They do not
address what happens when many agents discover the same interface repeatedly, who
maintains that knowledge as sites drift, or how the cost of discovery is shared --- the
maintenance-and-economics problem this paper's later layers take up.

\paragraph{Identity, transparency, and capabilities.} Our trust layer reuses
established cryptographic systems rather than inventing them: EdDSA
signatures~\cite{rfc8032}, zero-knowledge credential binding~\cite{zklogin,camlys},
Certificate-Transparency-style append-only logs~\cite{rfc6962}, the object-capability
model~\cite{ocap}, and the Dolev--Yao adversary~\cite{dolevyao}. The novelty is not the
primitives but their \emph{composition} into one identity that signs every layer and
one ledger that records every claim, with ERC-8004~\cite{erc8004} as the cross-agent
identity surface.

\paragraph{Agent payments and incentives.} x402 revives HTTP 402 as a real
micropayment rail~\cite{x402}. Prior agent-payment work largely stops at ``the agent
can pay.'' We extend it to ``the payment routes to everyone who created the value'' ---
indexer, domain owner, platform --- with the distribution and slashing analysis carried
by the Maintenance Network paper~\cite{fdry}, and we argue (Section on FDRY) that the
asset securing this must be fairly distributed for the trust to hold at all. To our
knowledge the combination --- signed multi-layer descent, a content-addressed-plus-
ledgered substrate, and a contribution-weighted three-party settlement over x402 --- is
not present in prior work as a single system.

\section{The stack as a self-similar tree}

We model the stack as a layered tree in the OSI tradition~\cite{osi}, but oriented
around \emph{who is acting} rather than \emph{which bytes move}:
\[
\text{screen-clicks} \to \text{browser} \to \text{CLI} \to \text{OS} \to
\text{kernel} \to \text{packet}.
\]
The end-to-end argument~\cite{e2e} supplies the rule for descent: a function lives
at the highest layer that can do it completely and correctly, and lower layers are
fallback, not first resort. A cached signed plan beats a CLI call beats a browser
action beats moving a mouse across pixels like it is 2004. But when the higher
layer simply cannot act --- no API, hostile anti-bot, a page that is 100\% JavaScript
and 0\% mercy --- the agent descends, and at the bottom it emits packets whose
TLS/HTTP fingerprint is indistinguishable from a real browser via
curl-impersonate~\cite{curlimpersonate}, matching exactly the JA3/JA4 signature
the other side is keying on~\cite{ja3,ja4}. \impl{} for browser/CLI/HTTP and
fingerprint-faithful fetch; \prop{} for a uniform \emph{signed} descent through the
OS and kernel layers.

The reason any of this composes is that the shape at every layer is identical:
observe, decide, sign, act. That self-similarity is the whole trick --- it is what
lets \emph{one} identity and \emph{one} cache serve the entire tree instead of six
bespoke integrations held together with hope and retry loops.

This is not a metaphor we are stretching. The whole stack is one instance of a
single structure --- the covenant tree from which our internal tooling is built ---
and every section below is one atom of it, mapped to a primitive that already
exists in the literature (Table~\ref{tab:atoms}). We did not invent these; we
aligned them.

\begin{table}[h]
\centering
\small
\begin{tabular}{@{}lll@{}}
\hline
\textbf{layer / atom} & \textbf{primitive} & \textbf{source} \\
\hline
descent (walk)       & curl-impersonate / JA3, JA4   & \cite{curlimpersonate,ja3,ja4} \\
network (economics)  & x402 split settlement / fair three-way compensation & \cite{x402,fdry} \\
agent record (node)  & ERC-8004 trustless agent      & \cite{erc8004} \\
the stack (tree)     & end-to-end argument / OSI     & \cite{e2e,osi} \\
witness (no leak)    & zkLogin / Pedersen commit     & \cite{zklogin,pedersen} \\
cache (value)        & Merkle content-addressing     & \cite{merkle} \\
\quad $\hookrightarrow$ KV parallel & attention KV cache / PagedAttention & \cite{attention,pagedattention} \\
seal (no unsigned)   & AEAD / object-capabilities    & \cite{rfc5116,ocap} \\
ledger (commitment)  & Certificate Transparency / hash chain & \cite{rfc6962,bitcoin} \\
settle (bonded)      & bonded route trust / FROST quorum & \cite{fdry,frost} \\
loop (control)       & OODA / proof-of-history clock & \cite{ooda,solana} \\
descent (walk)       & curl-impersonate / JA3, JA4   & \cite{curlimpersonate,ja3,ja4} \\
\hline
\end{tabular}
\caption{Each layer of the signed stack is one atom of the covenant tree, mapped
to an existing, cited primitive. A runnable reference implementation of the cache
+ ledger core ships alongside this paper (\texttt{reference/}, with tests that
execute each claim).}
\label{tab:atoms}
\end{table}

\section{Identity: one key signs every layer}
The root of trust is a single Ed25519 keypair~\cite{rfc8032} --- and conveniently,
the agent already has one, because a Solana account literally \emph{is} a base58
Ed25519 public key. The wallet is the identity. Every action --- a click, a
syscall, a packet --- is admissible only if it chains to a signature under that one
root. No signature, no action; the stack has no anonymous side door. \impl{}: the
Unbrowse signer already produces wallet signatures for resolve/execute admission.

For trust \emph{between} agents we adopt the ERC-8004 ``Trustless Agents''
model~\cite{erc8004}, which defines three on-chain registries: \textbf{Identity}
(portable agent IDs), \textbf{Reputation} (signed feedback), and
\textbf{Validation} (independent re-execution / ZK / TEE checks). An Unbrowse route
maintainer is precisely one of these agents: a portable identity that can carry
reputation and be checked by someone who does not trust it yet. \prop{}: the
ERC-8004 binding is a design target, not a shipped feature, and we will not pretend
otherwise three paragraphs from a contract address.

\section{Witness without disclosure: ZK credential binding}

Signing is the easy half. The hard half --- the actual research --- is proving a
credential belongs to the identity \emph{without revealing the credential}. Your
password, your session cookie, your 1Password entry, the fact that you clicked
``approve'' for this domain: all of it should be provably bound to the wallet and
leak \emph{nothing} identifiable at any layer unless you choose to open the proof
under signature.

This is the classical anonymous-credential problem~\cite{camlys}, and it has grown
up. zkLogin~\cite{zklogin} binds an existing OAuth/OpenID identity to a chain
address through a Groth16 proof, hiding the link even from the identity provider;
Semaphore~\cite{semaphore} proves anonymous group membership with zk-SNARKs;
Pedersen commitments~\cite{pedersen} give the ``commit now, open later, can't lie
about it'' primitive for a single secret. Together they make the binding a
two-witness corroboration that betrays nothing: the chain sees \emph{that} a
credential is bound, never \emph{what} it is. \prop{} --- this is the central new
proposal of the paper, and to be blunt, today Unbrowse stores credentials in a
boring local vault like everyone else. The proposal is the upgrade, not the status
quo.

\section{Cache: content-addressed, sealed unless revealed}

Every layer's settled result --- a resolved endpoint, a rendered page, a signed
plan --- is memoised under a content-addressed key, in the Merkle
tradition~\cite{merkle} and exactly as realised by content-addressed stores like
IPFS~\cite{ipfs} and Git's object model~\cite{gitobjects}. Same structure
re-walked, same key, same answer, deterministically; the cache is re-derived to
\emph{verify}, never trusted because it looked right last time. Unbrowse's
centralised cache servers can hold these entries and serve them across agents,
while a Pedersen commitment~\cite{pedersen} over a sensitive entry keeps it sealed
until the holder opens it --- a cache that is public by default and private by
proof. \impl{} for endpoint/route caching; \prop{} for the sealed-unless-revealed
commitment layer.


The parallel worth naming is to the \emph{KV cache} inside the transformer doing
the reasoning. Attention~\cite{attention} computes a key and value for every token;
the KV cache memoises those states so decoding never recomputes them, and modern
serving engines page that cache like OS virtual memory and \emph{share identical
prefixes across requests} by content (hash) key~\cite{pagedattention}. That is the
same discipline, one altitude up: same prefix, same cached blocks, reused not
recomputed --- except the structural key is now a covenant subtree rather than a
token prefix, and the entry is ZK-sealed rather than served in the clear. The model
caches computation it has already done; the stack caches \emph{actions} it has
already signed. Same shape, different altitude --- a KV cache all the way down,
where each layer's entry stays sealed until a signature opens it.
\section{Seal: nothing ships unsigned}

Before any layer emits, it self-verifies through the root. At the wire that is
authenticated encryption with associated data (AEAD)~\cite{rfc5116}: the receiver
checks the tag before it trusts a single byte, and tampering fails closed rather
than failing quietly into your logs. The same discipline generalises upward as
object-capability attenuation~\cite{ocap} --- authority travels only by reference
and can be narrowed, never picked up from the ambient air. The rule is boring on
purpose: no unsigned action crosses the gate, at any layer, ever. \prop{} as a
\emph{uniform} stack property --- AEAD itself is, obviously, ubiquitous and shipped
at the transport layer; we are not claiming to have invented encryption.

\section{Ledger: where the signatures actually land}

Signing is only half the story; the obvious next question --- and the one most
``signed'' systems wave away --- is \emph{where does the signature go}. A cache and
a ledger are not the same object, and conflating them is the bug. A cache is
content-addressed: you fetch a value by the hash of its content and order does not
matter, so the same hash resolves on any host~\cite{merkle}. A ledger is the
orthogonal half: append-only, \emph{ordered}, and hash-chained, so the history and
the order are themselves tamper-evident. A signed entry --- a signed KV result, a
route attestation, a bonded quality claim --- needs both: the cache holds the
value, the ledger holds the signed commitment to it.

The load-bearing design choice is \emph{value off-chain, root on-chain}. You never
put the payload on a chain --- that would be paying gas to store megabytes and
reading them slowly. You store the small, signed commitment: an entry of
\texttt{\{signer, value-hash, timestamp, prev-hash\}}, where the value itself lives
in the content-addressed cache and only its hash appears in the log. Each entry
chains to the prior, so reordering or editing any past row breaks every root after
it~\cite{bitcoin}. Independent auditors --- not a trusted operator --- verify the
log is append-only and consistent, exactly the model Certificate Transparency
formalises with its Merkle log, Signed Tree Head, and inclusion
proofs~\cite{rfc6962}. That is what makes the ledger trustless rather than
``trust our server.''

This also answers how the system decentralises without a rewrite. Per-entry
on-chain writes do not scale, so you batch: thousands of entries collapse to one
Merkle root, and only the \emph{root} is checkpointed on-chain (the
Certificate-Transparency / rollup pattern), ordered by a decentralised clock such
as Solana's proof of history~\cite{solana}. The maturity ladder is one shape at
every rung: a local append-only log today, a signed-root server next, on-chain
Merkle-root checkpoints as the decentralised end-state --- you swap the host, the
signatures and the hash-chain never change. \impl{} for the local append-only
ledger (the route graph already keeps signed JSONL records); \prop{} for the
on-chain root checkpointing and independent-auditor layer.

\section{From security to economics: where this paper stops}

A signed, sealed, ledgered route graph is a security substrate; what it is
\emph{worth} --- who maintains the routes, who pays, who earns, and how a freshness
claim is made costly to fake --- is a separate concern with its own paper. We draw
the line here deliberately. This paper is the security, authentication, and privacy
of the stack: one key signs every layer, credentials are bound by zero-knowledge
proof and revealed only under signature, every result is content-addressed and
sealed, and nothing ships unsigned. The maintenance economy that sits on top of
that substrate --- bonded accountability, proof of indexing, the challenge-and-slash
discipline, and the FDRY security stake --- is developed in full in the companion
Maintenance Network paper~\cite{fdry}.

Keeping the two concerns in separate papers is itself the discipline. A security
model must stand on its own, without leaning on a token to be true: the signatures,
the ZK binding, and the sealed cache are correct or not on cryptographic grounds
alone, independent of any economics layered above them. Conversely, the economics
is only worth stating once the substrate it secures is trustworthy. So we settle
the substrate here and hand the reader, deliberately, to the next paper for the
question of who keeps it fresh and how that claim is made accountable~\cite{fdry}.
\section{The control loop}

The whole stack runs as an OODA loop~\cite{ooda}: observe the layer's state, orient
it against the model (this is where the learning actually lives, despite ``decide''
getting all the press), decide the next action at the right layer, act under
signature --- then repeat with feedback, dropping a layer or replanning on failure,
and resting when the goal settles. It is the same plan/build/test/judge cycle, run
uniformly at every layer of the tree. If that sounds unglamorous, good: the
glamorous control loops are the ones that do not converge.



\section{Evaluation: the security substrate}

This paper's claims are about a security discipline, so we evaluate the security
substrate it introduces --- not the end-to-end latency win of route reuse, which is
a discovery and economics result developed in the companion paper~\cite{fdry}. We
are precise about what the measurements here do and do not establish.

\paragraph{The substrate is correct by executable test.} The local cache+ledger
reference implementation that ships with this paper is unit-verified:
content-addressing, value-off-chain/root-on-chain separation, hash-chain
tamper-evidence, ed25519 signatures, and deterministic Merkle-root inclusion all
pass as executable tests, so the substrate's correctness claims are runnable, not
asserted.

\paragraph{Sealed-cache reuse, measured.} We additionally ship a runnable
micro-benchmark (\texttt{reference/bench/\allowbreak bench\_reuse.py}) that isolates the one
property the sealed cache rests on --- content-addressed reuse versus re-derivation
--- on the actual cache implementation. Over 60 trials, re-deriving a representative
$\sim$64\,KB extracted payload averages \textbf{4.48\,ms} while a content-addressed
cache hit (read + hash re-verify) averages \textbf{0.052\,ms}: a measured
\textbf{$\approx$95$\times$} mean speedup, wall-clock, deterministic, with a gate
that fails if reuse is not materially faster. This is a CPU-only
recompute-versus-reuse demonstration --- it isolates the cache primitive, not the
network or the browser --- and it establishes exactly one thing: a sealed,
content-addressed commitment is asymptotically cheaper to re-read than to
re-derive, which is what makes the cache safe to lean on rather than merely fast.

\paragraph{What this establishes, and what it does not.} The substrate tests
establish that the security primitives are correct as implemented; the
micro-benchmark establishes that sealed-cache reuse is cheap. Neither establishes
end-to-end task success, cold-start discovery cost, or accuracy on adversarial,
heavily-JavaScript-gated sites where the descent to the browser layer is forced ---
those are out of this paper's scope and are taken up, where they belong, in the
companion Maintenance Network paper~\cite{fdry}.
\section{Threat model}

We state the adversary explicitly, because a trust paper that never names what it
defends against is decoration. We assume a Dolev--Yao network
attacker~\cite{dolevyao} --- able to read, drop, replay, and forge messages on any
layer --- plus three application-level adversaries the route economy specifically
invites. For each, we name the atom that resists it and the residual risk we do not
claim to close.

\paragraph{(A1) The impersonator.} An attacker replays or forges an action to act
as another identity. Resisted at the \emph{root}: every action carries an Ed25519
signature over its canonical bytes~\cite{rfc8032}, and the ledger's
\texttt{prev}-hash makes a replayed entry detectable out of order~\cite{rfc6962}.
Residual: key theft is out of scope --- a stolen wallet key is the user, by
construction. This is the same boundary every signing system accepts.

\paragraph{(A2) The credential thief / over-reach.} An operator retargets imported
sessions at accounts they do not own --- the abuse vector that closed the
source~\cite{ossnotice}. The \emph{witness} atom narrows it: credentials are bound
to the identity by zero-knowledge proof and revealed only under
signature~\cite{zklogin,camlys}, so a leaked ledger or cache exposes \emph{that} a
credential is bound, never the credential. Capability attenuation~\cite{ocap} bounds
what a delegated action may do. Residual: an operator authenticated as themselves
can still drive their own credentials; we constrain reach, not self-harm.

\paragraph{(A3) The poisoner.} An adversary publishes a false or stale route to
mislead later agents --- the integrity attack on a shared graph. Resisted at
\emph{settle} and \emph{witness} together: a trust-tier route carries a slashable
bond and is challengeable on independent evidence~\cite{fdry}, and finalisation
needs a $t$-of-$n$ quorum~\cite{frost}, so one dishonest maintainer cannot settle a
claim. Residual: open (unbonded) routes remain best-effort by design; the model
moves trust-sensitive traffic to bonded routes, it does not claim the open tier is
adversary-proof.

\paragraph{(A4) The free-rider and the Sybil.} An actor consumes the graph without
paying, or splits into many identities to farm rewards. The \emph{verb} atom binds
payment to the act: an unpaid \texttt{execute} gets HTTP 402, not
service~\cite{x402}. Sybil farming is bounded because reward shares derive from
on-ledger contribution metrics, not headcount --- minting identities does not mint
contribution. Residual: collusion between an indexer and a domain owner to inflate a
route's apparent value is a real open problem; we flag it rather than wave it away,
and it is the strongest argument for the bonded tier's slashing teeth.

\paragraph{(A5) The exfiltrator of the moat.} An adversary --- or an honest
contributor by accident --- leaks the closed capture/integrity engine into a public
artifact. Resisted mechanically, not by policy: \texttt{leak-guard.sh} and
\texttt{paper-gate.sh} run in release CI and fail the build if a sensitive term or
an unanchored claim reaches a public path. This paper is itself subject to that
gate. Residual: NDA source review remains the only full-disclosure path, by the
abuse reasoning of~\cite{ossnotice}.

\paragraph{What we do not defend.} We make no claim against a malicious endpoint
that lies in its own responses (a route can faithfully replay a hostile API), a
global passive adversary correlating timing across the whole network, or coercion of
a key holder. These are out of scope and named so the reader is not misled by
silence.

\section{What is actually built (no fabricated green)}

In the spirit of not selling a roadmap as a changelog:
\begin{itemize}[leftmargin=1.4em]
\item \impl{} Intent $\to$ route resolve $\to$ execute; live browser capture; HTTP
fetch with browser-faithful TLS fingerprinting; wallet-signed admission;
route/endpoint caching.
\item \prop{} Uniform \emph{signed} descent through the OS/kernel layers; ZK
credential binding (the central contribution); the sealed-unless-revealed
commitment cache. The bonded route-trust economy (ERC-8004 binding, FDRY bonding /
challenge / slashing, and proof of indexing) is the subject of the companion
Maintenance Network paper~\cite{fdry}, not claimed here.
\end{itemize}
Treat the \prop{} list as a research agenda with cited foundations, not a feature
list with a release date. The gap between those two things is the entire
difference between a whitepaper and a press release, and we would like to stay on
the correct side of it.

\section{Conclusion}

Internal APIs were an excellent first layer --- enough to prove the wedge in a real
market, which is more than most ideas manage. They were never going to be the last
layer, because an agent that is sovereign over its own actions has to be one
coherent entity across every layer it touches: under one key, revealing nothing it
did not choose to reveal, and standing behind its routes with something it can
actually lose. The contribution is deliberately modest --- one covenant (sign it,
seal it, cache it, bind identity without disclosure) that holds all the way down,
each layer pinned to a real primitive. Attention, it turns out, was a great deal of
what you needed. It just was never the \emph{stack}.

\begin{thebibliography}{99}
\bibitem{ossnotice} Unbrowse AI. \emph{Open Source Notice --- the closed/open boundary}. docs/OPEN-SOURCE-NOTICE.md, 2026.
\bibitem{x402} Coinbase. \emph{x402: An Open Protocol for Internet-Native
Payments over HTTP 402}. x402.org whitepaper, 2025. \url{https://github.com/coinbase/x402}.
\bibitem{fdry} L.~Tham. \emph{Unbrowse Maintenance Network: Proof of Indexing and
Bonded Accountability in a Shared Route Graph}. Unbrowse AI, 2026.
\bibitem{e2e} J.~H.~Saltzer, D.~P.~Reed, D.~D.~Clark. \emph{End-to-End Arguments
in System Design}. ACM TOCS 2(4):277--288, 1984. DOI: 10.1145/357401.357402.
\bibitem{osi} ITU-T Recommendation X.200 (= ISO/IEC 7498-1), \emph{OSI Basic
Reference Model}, 1994.
\bibitem{react} S.~Yao, J.~Zhao, et al. \emph{ReAct: Synergizing Reasoning and Acting in Language Models}. ICLR 2023. arXiv:2210.03629.
\bibitem{mind2web} X.~Deng, Y.~Gu, et al. \emph{Mind2Web: Towards a Generalist Agent for the Web}. NeurIPS 2023. arXiv:2306.06070.
\bibitem{webarena} S.~Zhou, F.~F.~Xu, et al. \emph{WebArena: A Realistic Web Environment for Building Autonomous Agents}. ICLR 2024. arXiv:2307.13854.
\bibitem{seeact} B.~Zheng, B.~Gou, et al. \emph{GPT-4V(ision) is a Generalist Web Agent, if Grounded (SeeAct)}. ICML 2024. arXiv:2401.01614.
\bibitem{toolformer} T.~Schick, J.~Dwivedi-Yu, et al. \emph{Toolformer: Language Models Can Teach Themselves to Use Tools}. NeurIPS 2023. arXiv:2302.04761.
\bibitem{gorilla} S.~G.~Patil, T.~Zhang, et al. \emph{Gorilla: Large Language Model Connected with Massive APIs}. NeurIPS 2024. arXiv:2305.15334.
\bibitem{mcp} Anthropic. \emph{Model Context Protocol}. 2024. \url{https://modelcontextprotocol.io}.
\bibitem{dolevyao} D.~Dolev, A.~C.~Yao. \emph{On the Security of Public Key Protocols}. IEEE Trans. Information Theory 29(2), 1983. DOI: 10.1109/TIT.1983.1056650.
\bibitem{rfc8032} S.~Josefsson, I.~Liusvaara. \emph{Edwards-Curve Digital
Signature Algorithm (EdDSA)}. RFC 8032, IRTF CFRG, 2017.
\bibitem{erc8004} \emph{ERC-8004: Trustless Agents}. Ethereum Improvement
Proposals (Draft). \url{https://eips.ethereum.org/EIPS/eip-8004}
\bibitem{camlys} J.~Camenisch, A.~Lysyanskaya. \emph{An Efficient System for
Non-transferable Anonymous Credentials with Optional Anonymity Revocation}.
EUROCRYPT 2001, LNCS 2045. IACR ePrint 2001/019.
\bibitem{zklogin} F.~Baldimtsi, K.~Chalkias, et al. \emph{zkLogin:
Privacy-Preserving Blockchain Authentication with Existing Credentials}. ACM CCS
2024. arXiv:2401.11735.
\bibitem{semaphore} Semaphore Protocol. \emph{Semaphore: anonymous signaling on
Ethereum}. \url{https://github.com/semaphore-protocol/semaphore}
\bibitem{pedersen} T.~P.~Pedersen. \emph{Non-Interactive and
Information-Theoretic Secure Verifiable Secret Sharing}. CRYPTO 1991, LNCS 576.
DOI: 10.1007/3-540-46766-1\_9.
\bibitem{merkle} R.~C.~Merkle. \emph{A Digital Signature Based on a Conventional
Encryption Function}. CRYPTO 1987, LNCS 293. DOI: 10.1007/3-540-48184-2\_32.
\bibitem{rfc6962} B.~Laurie, A.~Langley, E.~Kasper. \emph{Certificate
Transparency}. RFC 6962, IETF, 2013.
\bibitem{bitcoin} S.~Nakamoto. \emph{Bitcoin: A Peer-to-Peer Electronic Cash
System}. 2008. \url{https://bitcoin.org/bitcoin.pdf}
\bibitem{solana} A.~Yakovenko. \emph{Solana: A New Architecture for a High
Performance Blockchain}. 2018. \url{https://solana.com/solana-whitepaper.pdf}
\bibitem{ipfs} Protocol Labs. \emph{IPFS / IPLD Merkle-DAG specifications}.
\bibitem{attention} A.~Vaswani, N.~Shazeer, et al. \emph{Attention Is All You
Need}. NeurIPS 2017. arXiv:1706.03762.
\bibitem{pagedattention} W.~Kwon, Z.~Li, et al. \emph{Efficient Memory Management
for Large Language Model Serving with PagedAttention}. SOSP 2023. arXiv:2309.06180;
\url{https://github.com/vllm-project/vllm}.
\url{https://github.com/ipfs/specs}
\bibitem{gitobjects} S.~Chacon, B.~Straub. \emph{Pro Git}, 2nd ed., \S10.2 Git
Internals --- Git Objects. \url{https://git-scm.com/book/en/v2/Git-Internals-Git-Objects}
\bibitem{rfc5116} D.~McGrew. \emph{An Interface and Algorithms for Authenticated
Encryption}. RFC 5116, 2008.
\bibitem{ocap} M.~S.~Miller. \emph{Robust Composition: Towards a Unified Approach
to Access Control and Concurrency Control}. PhD thesis, Johns Hopkins University,
2006.
\bibitem{curlimpersonate} curl-impersonate. \emph{A special build of curl that
impersonates real browsers (TLS/HTTP fingerprints)}.
\url{https://github.com/lwthiker/curl-impersonate}
\bibitem{ja3} Salesforce. \emph{JA3: TLS client fingerprinting}.
\url{https://github.com/salesforce/ja3}
\bibitem{ja4} FoxIO. \emph{JA4+ network fingerprinting suite}.
\url{https://github.com/FoxIO-LLC/ja4}
\bibitem{frost} C.~Komlo, I.~Goldberg. \emph{FROST: Flexible Round-Optimized
Schnorr Threshold Signatures}. SAC 2020, LNCS 12804. IACR ePrint 2020/852; RFC 9591.
\bibitem{ooda} J.~R.~Boyd. \emph{The Essence of Winning and Losing}, 1995; in
\emph{A Discourse on Winning and Losing}, ed. G.~Hammond, Air University Press, 2018.
\end{thebibliography}

\end{document}
